\chapter{Complexity — Clean}
%%% generated by the blueprint pipeline from TCSlib.Complexity.SpaceComplexity.Machines.Clean
\section{Overview}

Declarations of the Lean module \texttt{TCSlib.Complexity.SpaceComplexity.Machines.Clean}.

\section{Declarations}

\begin{definition}[Marking a distinguished set of integers]
\lean{Complexity.LogProg.markSet}
\label{Complexity.LogProg.markSet}
\leanok
For a finite set $F$ of integers, define a partial Boolean marking on the integers by
assigning true at $0$, false at every member of $F$, and no value at all other integers.
\end{definition}

\begin{definition}[Simulation state assignment]
\lean{Complexity.LogProg.simSt}
\label{Complexity.LogProg.simSt}
\leanok
\uses{Complexity.LogProg.CleanSt, Turing.Cfg}
For a configuration of a $k$-work-tape Turing machine with Boolean input, assign a
simulation state as follows: if the machine has internal state $q$, the assigned state
is the simulation state corresponding to $q$. If the machine is halted and $k>0$, assign
the cleaning state for tape $0$ in the marking phase; if it is halted and $k=0$, assign
no state.
\end{definition}

\begin{definition}[Visited work-tape head positions]
\lean{Complexity.LogProg.visB}
\label{Complexity.LogProg.visB}
\leanok
\uses{Turing.Cfg.init, Turing.MultiTapeTM.runFrom}
Let $M$ be a $k$-work-tape Turing machine with Boolean input alphabet, let $q$ be an
initial state, and let $V$ be a Boolean input string. For a natural number $t$ and a
tape $i$, the visited-position set consists of the positions of tape $i$'s head in the
configurations reached from this initial setup at times $0,1,\ldots,t-1$; in particular,
it is empty when $t=0$.
\end{definition}

\begin{definition}[Simulated machine configuration]
\lean{Complexity.LogProg.simCfg}
\label{Complexity.LogProg.simCfg}
\leanok
\uses{Complexity.LogProg.CleanSt, Complexity.LogProg.ccfg, Complexity.LogProg.markSet, Complexity.LogProg.simSt, Complexity.LogProg.visB, Turing.Cfg, Turing.Cfg.init, Turing.MultiTapeTM.runFrom}
Let $M$ be a $k$-work-tape Turing machine with Boolean input alphabet, let $q$ be an
initial state, let $V$ be a Boolean input string, and let $t$ be a natural number. The
associated configuration of a machine with $2k$ work tapes has the simulation state
assigned to the original machine's configuration at time $t$, and uses that
configuration's input-head position, work-tape contents and head positions, and output.
Its additional work-tape data is formed from the visited-position sets for the original
run through times $0$ to $t-1$, together with the original head positions.
\end{definition}

\begin{lemma}[Initial cleaning-machine step]
\lean{Complexity.LogProg.init_step}
\label{Complexity.LogProg.init_step}
\leanok
\uses{Complexity.LogProg.ccfg, Complexity.LogProg.cleanTM, Complexity.LogProg.idleK, Complexity.LogProg.markSet, Complexity.LogProg.simCfg, Complexity.LogProg.visB, Turing.Action.apply, Turing.Cfg.init, Turing.MultiTapeTM.runFrom_zero, Turing.MultiTapeTM.step, Turing.moveInputPos_zero}
\difficulty{4}
Let $M$ be a machine with $k$ work tapes, and let $q_1$ and $q$ be states and $V$ an
input word. One step of the cleaning machine constructed from $M$ with initial state
$q_1$, starting from its initial configuration with state $\mathrm{init}(q)$ on $V$,
yields the simulation configuration for $M$, started in state $q$ on $V$, at time zero.
\end{lemma}

\begin{lemma}[Successive visited-position sets]
\lean{Complexity.LogProg.visB_succ}
\label{Complexity.LogProg.visB_succ}
\leanok
\uses{Complexity.LogProg.visB, Turing.Cfg.init, Turing.MultiTapeTM.runFrom}
\difficulty{1}
For a multi-tape Turing machine $M$, initial state $q$, input word $V$, time
$t\in\mathbb N$, and work tape $i$, the set of head positions visited from time $0$
through time $t+1$ is the set of positions visited from time $0$ through time $t$, with
the head position at time $t$ added.
\end{lemma}

\begin{lemma}[Writing the marked set]
\lean{Complexity.LogProg.wr_markSet}
\label{Complexity.LogProg.wr_markSet}
\leanok
\uses{Complexity.LogProg.markSet, Complexity.LogProg.wr}
\difficulty{5}
For any finite set $F$ of integers and integer $p$, start with the partial Boolean
marking that assigns true at $0$, false on $F$, and is undefined elsewhere. If its value
at $p$ is true, leave it unchanged; otherwise write false at $p$. The resulting marking
is the one that assigns true at $0$, false on $F\cup\{p\}$, and is undefined elsewhere.
\end{lemma}

\begin{lemma}[Simulation Step with Cleanup Tracks]
\lean{Complexity.LogProg.step_sim}
\label{Complexity.LogProg.step_sim}
\leanok
\uses{Complexity.LogProg.CleanSt, Complexity.LogProg.ccfg, Complexity.LogProg.ccfg_mark, Complexity.LogProg.ccfg_state, Complexity.LogProg.ccfg_work, Complexity.LogProg.cleanTM, Complexity.LogProg.simSt, Complexity.LogProg.wr, Turing.Action.apply, Turing.Cfg, Turing.Cfg.inputSymbol, Turing.Cfg.workTapeSymbols, Turing.MultiTapeTM, Turing.MultiTapeTM.step}
\difficulty{5}
Let $M$ be a Boolean-alphabet Turing machine with $k$ work tapes, and let $c$ be a
configuration whose state is $q'$. For any contents $M_t$ of an additional set of $k$
work tapes, one step of the cleaning machine started in simulation state $q'$ with the
original tapes and head positions from $c$, and with the additional tapes having
contents $M_t$ and head positions matching those of $c$, yields the configuration whose
simulated state and original tapes and output are those of one step of $M$ from $c$. The
additional tapes retain their head positions, and each is changed only at the old head
position: there it is left unchanged if its content was $1$, and otherwise is set to $1$
when it was already $1$ or to $0$ when it was not; equivalently, the scanned cell is
unchanged when it contains $1$ and is set to $0$ otherwise.
\end{lemma}

\begin{lemma}[One-step simulation of a live computation]
\lean{Complexity.LogProg.sim_step_clean}
\label{Complexity.LogProg.sim_step_clean}
\leanok
\uses{Complexity.LogProg.ccfg, Complexity.LogProg.cleanTM, Complexity.LogProg.markSet, Complexity.LogProg.simCfg, Complexity.LogProg.simSt, Complexity.LogProg.step_sim, Complexity.LogProg.visB, Complexity.LogProg.visB_succ, Complexity.LogProg.wr_markSet, Turing.Cfg.init, Turing.MultiTapeTM, Turing.MultiTapeTM.runFrom, Turing.MultiTapeTM.runFrom_succ_eq_step', Turing.MultiTapeTM.step}
\difficulty{4}
Let $M$ be a deterministic machine with $k$ Boolean work tapes, let $q_1,q$ be states,
let $V$ be a Boolean input, and let $t\ge 0$. If the configuration reached by $M$ from
state $q$ on input $V$ after $t$ steps is not halted, then one step of the cleaning
machine initialized with state $q_1$, from the configuration encoding that simulated
configuration, equals the encoded configuration reached by $M$ after $t+1$ steps.
\end{lemma}

\begin{lemma}[Visited positions form an interval]
\lean{Complexity.LogProg.visited_interval}
\label{Complexity.LogProg.visited_interval}
\leanok
\uses{Turing.Cfg.init, Turing.MultiTapeTM, Turing.MultiTapeTM.ConfigCount.exists_eq_of_between, Turing.MultiTapeTM.runFrom, Turing.MultiTapeTM.runFrom_succ_eq_step', Turing.MultiTapeTM.visitedByTapeHead, Turing.MultiTapeTM.workTapePos_step_le}
\difficulty{6}
For every multi-tape Turing machine $M$, initial state $q$, Boolean input $V$, time
$T\in\mathbb N$, and work tape $j$, the head positions visited from time $0$ through
time $T$ form an integer interval $[a,b]$ for some integers $a\le 0\le b$, and both
endpoints $a$ and $b$ are visited.
\end{lemma}

\begin{definition}[Cleaning-phase state assignment]
\lean{Complexity.LogProg.stSt}
\label{Complexity.LogProg.stSt}
\leanok
\uses{Complexity.LogProg.CleanSt}
Let $k$ be a natural number and let $S$ be a state type. For a natural number $i$,
assign the cleaning state for tape $i$ in the marking phase when $i<k$; when $i\ge k$,
assign no state.
\end{definition}

\begin{definition}[Cleaning-stage configuration]
\lean{Complexity.LogProg.stg}
\label{Complexity.LogProg.stg}
\leanok
\uses{Complexity.LogProg.CleanSt, Complexity.LogProg.ccfg, Complexity.LogProg.stSt, Turing.Cfg}
Let $k$ be a natural number, let $S$ be a state type, and let $x$ be a Boolean input
string. Given a natural number $i$, an input-head position among the $|x|+2$ allowed
positions, work-tape contents $D_t,D_m$ on the integers, work-tape head positions $D_p$,
and a Boolean output string, form a configuration with $2k$ work tapes. Its state is the
cleaning state for tape $i$ in the marking phase if $i<k$, and is absent otherwise. In
each group of $k$ tapes, tapes with index less than $i$ are blank with head at position
zero; the other tapes use the corresponding contents and head positions from
$D_t,D_m,D_p$.
\end{definition}

\begin{lemma}[Bounded cleaning stages]
\lean{Complexity.LogProg.stages_run}
\label{Complexity.LogProg.stages_run}
\leanok
\uses{Complexity.LogProg.CleanSt, Complexity.LogProg.PB, Complexity.LogProg.PB_ccfg, Complexity.LogProg.ccfg, Complexity.LogProg.clNext, Complexity.LogProg.cleanTM, Complexity.LogProg.clean_tape, Complexity.LogProg.markI, Complexity.LogProg.runFrom_forall_append, Complexity.LogProg.stSt, Complexity.LogProg.stg, Complexity.LogProg.wr, Turing.Cfg, Turing.MultiTapeTM, Turing.MultiTapeTM.runFrom, Turing.MultiTapeTM.runFrom_add}
\difficulty{6}
Let $M$ be a deterministic machine with $k$ Boolean work tapes, let $q_1$ be a state,
and fix an input, input-head position, tape contents $D_t,M_t$, head positions $D_p$,
output word, and integer $B\ge 0$. Assume that for each tape $j$ there are integers
$a,b$ with $a\le0\le b$, $-B\le a-1$, $b+1\le B$, $a\le D_p(j)\le b$, and such that
changing the marked-tape contents at $D_p(j)$ according to the prescribed write rule
produces the tape that marks the origin and every cell from $a$ through $b$; also assume
every nonblank cell of $D_t(j)$ lies between $a$ and $b$. Then, for every $n,i\in\mathbb
N$ with $i+n=k$, there is a finite time $T$ such that throughout the cleaning-machine
run from stage $i$, all work-tape head positions have absolute value at most $B$, and at
time $T$ the configuration is stage $k$ with the same input-head position, tape data,
head positions, and output.
\end{lemma}

\begin{theorem}[Clean Simulation with Tape Erasure]
\lean{Complexity.LogProg.cleanTM_run}
\label{Complexity.LogProg.cleanTM_run}
\leanok
\uses{Complexity.LogProg.CleanSt, Complexity.LogProg.PB, Complexity.LogProg.PB_ccfg, Complexity.LogProg.ccfg, Complexity.LogProg.cleanTM, Complexity.LogProg.init_step, Complexity.LogProg.markI, Complexity.LogProg.markSet, Complexity.LogProg.runFrom_forall_append, Complexity.LogProg.simCfg, Complexity.LogProg.simSt, Complexity.LogProg.sim_step_clean, Complexity.LogProg.stSt, Complexity.LogProg.stages_run, Complexity.LogProg.stg, Complexity.LogProg.visB, Complexity.LogProg.visited_interval, Complexity.LogProg.wr, Complexity.LogProg.wr_markSet, Turing.Cfg, Turing.Cfg.init, Turing.MultiTapeTM, Turing.MultiTapeTM.abs_pos_lt_card_visited, Turing.MultiTapeTM.initCfg, Turing.MultiTapeTM.mem_visited_of_ne_none, Turing.MultiTapeTM.runFrom, Turing.MultiTapeTM.runFrom_add, Turing.MultiTapeTM.runFrom_of_halt, Turing.MultiTapeTM.runFrom_succ_eq_step', Turing.MultiTapeTM.visitedByTapeHead, Turing.MultiTapeTM.visitedByTapeHead_mono}
\difficulty{6}
Let $M$ be a multi-tape Turing machine with $k$ Boolean work tapes and state set $S$,
and choose states $q_1,q\in S$, a Boolean input word $V$, a time $T$, and a natural
number $s$. Assume that $M$, started in state $q$ on $V$, has halted at time $T$, and
that by time $T$ the set of cells visited by the head of each work tape has cardinality
at most $s$. Then the cleaning machine started in its initialization state carrying $q$
on input $V$ eventually halts with the same output as $M$ at time $T$, with all its
work-tape cells blank and all its work-tape heads at the origin; throughout this run,
every work-tape head has absolute position at most $s$.
\end{theorem}

\begin{lemma}[Monotonicity of the Work-Tape Bound]
\lean{Complexity.LogProg.CleanRun.mono}
\label{Complexity.LogProg.CleanRun.mono}
\leanok
\uses{Complexity.LogProg.CleanRun, Turing.MultiTapeTM}
\difficulty{1}
Let $D$ be a multitape Turing machine, with starting state $q$, Boolean input word $V$,
answer bit $b$, and natural-number bounds $B$ and $B'$. If $D$ has a run that halts with
output exactly $[b]$, leaves every work-tape cell blank and every work-tape head at
position zero, and keeps every work-tape head within absolute position $B$ at every time
up to and including its halting time, then whenever $B\le B'$, it also has such a run
with the bound $B'$.
\end{lemma}
